\documentclass[11pt]{amsart}
\usepackage[localtheorems]{raeez-math-template}

\providecommand{\C}{\mathbb{C}}
\providecommand{\dbar}{\bar{\partial}}
\providecommand{\g}{\mathfrak{g}}
\providecommand{\Sym}{\mathrm{Sym}}
\providecommand{\CE}{\mathrm{CE}}
\providecommand{\chir}{\mathrm{chir}}
\providecommand{\Ran}{\mathrm{Ran}}
\providecommand{\Hom}{\mathrm{Hom}}
\providecommand{\cL}{\mathcal{L}}
\providecommand{\cM}{\mathcal{M}}
\providecommand{\cD}{\mathcal{D}}
\providecommand{\cO}{\mathcal{O}}
\providecommand{\cT}{\mathcal{T}}
\providecommand{\Linf}{L_\infty}
\providecommand{\Jet}{J^\infty}

\theoremstyle{plain}
\newtheorem{theorem}{Theorem}[section]
\newtheorem{proposition}[theorem]{Proposition}
\newtheorem{lemma}[theorem]{Lemma}
\theoremstyle{definition}
\newtheorem{definition}[theorem]{Definition}
\newtheorem{construction}[theorem]{Construction}
\theoremstyle{remark}
\newtheorem{remark}[theorem]{Remark}

\title{Chiral Chevalley--Eilenberg Cohomology of a Holomorphic
$\Linf$-Algebroid}
\date{}

\begin{document}
\maketitle

\section{Holomorphic algebroids}

Let $X$ be a complex manifold. A holomorphic $\Linf$-algebroid is the
Dolbeault resolution of a holomorphic algebroid together with higher
brackets. It supplies the infinitesimal symmetries whose chiral
Chevalley--Eilenberg complex records local observables on the Ran space
$\Ran X$.

\begin{definition}[Holomorphic $\Linf$-algebroid]\label{def:hol-linf-algebroid}
A holomorphic $\Linf$-algebroid on $X$ is a graded
$\Omega_X^{0,\bullet}$-module $\cL$ with a degree-one differential
$l_1$, a holomorphic anchor
\[
  \rho \colon \cL \longrightarrow \Omega_X^{0,\bullet}(\cT_X),
\]
and multilinear brackets
\[
  l_n \colon \Lambda^n_{\Omega_X^{0,\bullet}}\cL
  \longrightarrow \cL[2-n],
  \qquad n\geq 2,
\]
satisfying the $\Linf$ Jacobi identities. The binary bracket satisfies
the anchored Leibniz rule
\[
  l_2(x, f y)
  =
  \rho(x)(f)y + (-1)^{|x||f|}f\,l_2(x,y)
\]
for local sections $x,y$ and $f\in\Omega_X^{0,\bullet}$, and the higher
brackets are $\Omega_X^{0,\bullet}$-multilinear after the usual
algebroid correction in the anchored slot.
\end{definition}

\begin{definition}[Chiral module]\label{def:chiral-module}
A chiral $\cL$-module is a Dolbeault complex $\cM$ equipped with
operations
\[
  a_n \colon \Lambda^{n-1}\cL\otimes\cM\longrightarrow \cM[2-n],
  \qquad n\geq 1,
\]
compatible with the brackets of $\cL$, the anchor, and the small-diagonal
support condition on $\Ran X$. The compatibility identities are the
module part of the $\Linf$ Jacobi system.
\end{definition}

Beilinson--Drinfeld chiral Lie algebras on curves are encoded by
right $\cD$-modules and brackets supported on diagonals. In the present
holomorphic setting the $\cD$-module direction is the $(1,0)$ direction,
while the Dolbeault differential resolves the antiholomorphic direction.
Francis--Gaitsgory's Ran-space formalism supplies the factorization
envelope in arbitrary complex dimension once the operations are defined
on infinitesimal neighborhoods of the diagonals.

\section{The chiral Chevalley--Eilenberg complex}

\begin{definition}[Chiral Chevalley--Eilenberg complex]
\label{def:chiral-ce-complex}
For a holomorphic $\Linf$-algebroid $\cL$ and a chiral $\cL$-module
$\cM$, define the completed local Chevalley--Eilenberg complex by
\[
  \CE_{\dbar,\chir}^{p,q}(\cL,\cM)
  :=
  \Hom_{\Omega_X^{0,\bullet}}
  \left(
    \widehat{\Sym}^{p}\bigl(\cL^\vee[1]\bigr),
    \Omega_X^{0,q}(\cM)
  \right)^{\Sigma_p}.
\]
The completion is taken in the symmetric degree and in the order of jets
along the small diagonals of $\Ran X$. The total differential is
\[
  D = d_{\CE}+\dbar_{\cM},
\]
where $d_{\CE}$ is the coderivation determined by the transposed
brackets $l_n^\vee$ and by the module operations $a_n$.
\end{definition}

\begin{proposition}[Differential]\label{prop:D-square-zero}
If the brackets, anchor, and module operations form a holomorphic
$\Linf$-algebroid-module pair, then
\[
  D^2=0.
\]
\end{proposition}

\begin{proof}
The square of $d_{\CE}$ is the Chevalley--Eilenberg expression for the
$\Linf$ Jacobi identities and the module identities. The Dolbeault term
satisfies $\dbar_{\cM}^2=0$. Holomorphicity of the anchor and of the
brackets gives
\[
  [d_{\CE},\dbar_{\cM}]=0.
\]
Thus
\[
  D^2=d_{\CE}^2+[d_{\CE},\dbar_{\cM}]+\dbar_{\cM}^2=0.
\]
\end{proof}

\begin{remark}[Curved inputs]\label{rem:curved-inputs}
For a curved holomorphic $\Linf$-algebroid the curvature
$l_0\in\cL[-1]$ contributes the usual contraction term to $d_{\CE}$.
The same proof gives $D^2=0$ precisely when the Maurer--Cartan equation
for $l_0,l_1,l_2,\ldots$ holds. Without that equation the square is
contraction by the curvature defect.
\end{remark}

\section{Jets and Ran locality}

\begin{construction}[Jet algebroid]\label{constr:jet-algebroid}
Let $\Jet\cL$ be the pro-holomorphic jet bundle of $\cL$. Its sections
are compatible formal Taylor expansions of local sections of $\cL$.
The brackets $l_n$ and the anchor lift functorially to $\Jet\cL$, and
the lifted operations are polynomial in the jet variables. The same
construction sends $\cM$ to $\Jet\cM$.
\end{construction}

\begin{proposition}[Ran-space factorization]\label{prop:ran-locality}
Assume that $\cL$ and $\cM$ are complete for the jet filtration and that
the horizontal $(1,0)$ connection on jets commutes with the Dolbeault
differential. Then
\[
  \CE_{\dbar,\chir}^{*}(\cL,\cM)
  \simeq
  \CE_{\dbar,\chir}^{*}(\Jet\cL,\Jet\cM)^{\mathrm{hor}}
\]
as filtered complexes. The right-hand side is a factorization algebra on
$\Ran X$.
\end{proposition}

\begin{proof}
The horizontal jet resolution is a filtered resolution of holomorphic
sections. The associated graded complex is the completed symmetric
algebra on holomorphic jets, and the horizontal differential removes the
formal Taylor variables. The Dolbeault differential is transverse to the
horizontal $(1,0)$ connection, so the filtered bicomplex has commuting
horizontal and Dolbeault differentials. The diagonal support of the
chiral operations is preserved by jet completion, giving the Ran-space
factorization product.
\end{proof}

\begin{remark}[Comparison data]\label{rem:comparison-data}
A compact or equivariant factorization comparison uses more data than
the formal local construction: an orientation of the determinant line,
a framing of the CY trace, a completion compatible with the monoid of
effective classes, equivariance for the symmetry group being used, and
a comparison map preserving the brackets and filtrations. The proposition
constructs the local chiral CE complex before these global comparison
choices are imposed.
\end{remark}

\section{\texorpdfstring{$K3\times E$}{K3xE} current algebra}

Let $X=K3\times E$, where $E$ is an elliptic curve, and let $\g$ be a
finite-dimensional complex Lie algebra. The Dolbeault current
algebroid is
\[
  \cL_X:=\Omega^{0,\bullet}(X,\g)
  =
  \g\otimes_{\C}\Omega^{0,\bullet}(X),
\]
with differential $\dbar$, zero anchor, binary bracket
\[
  [a\otimes\alpha,b\otimes\beta]
  =
  [a,b]_{\g}\otimes\alpha\wedge\beta,
\]
and $l_n=0$ for $n\geq3$. The coefficient module
$\cM=\Omega^{0,\bullet}(X)$ carries the trivial current action.

\begin{proposition}[Current CE cohomology]\label{prop:k3e-current-ce}
There is a filtered quasi-isomorphism
\[
  \CE_{\dbar,\chir}^{*}(\cL_X,\Omega^{0,\bullet}(X))
  \simeq
  \CE_{\chir}^{*}
  \left(\g\otimes_{\C}H^{0,\bullet}(X),H^{0,\bullet}(X)\right)
\]
where the right-hand side uses the cup-product algebra
$H^{0,\bullet}(X)$ and the induced current bracket. The Euler
supertrace of the Dolbeault factor is
\[
  \sum_q(-1)^q h^{0,q}(K3\times E)=0.
\]
\end{proposition}

\begin{proof}
Hodge theory gives a contraction of the Dolbeault dga
$\Omega^{0,\bullet}(X)$ onto the harmonic algebra $H^{0,\bullet}(X)$.
Tensoring this contraction with the finite-dimensional Lie algebra
$\g$ transfers the dg Lie current algebra to
$\g\otimes_\C H^{0,\bullet}(X)$ by the homological perturbation lemma.
The completed Chevalley--Eilenberg functor sends this filtered
$\Linf$ quasi-isomorphism to a filtered quasi-isomorphism of completed
CE complexes. The Hodge numbers factor by K\"unneth:
\[
  \chi(\cO_{K3\times E})
  =
  \chi(\cO_{K3})\chi(\cO_E)
  =
  2\cdot 0
  =
  0.
\]
\end{proof}

\begin{proposition}[Separated invariants on \texorpdfstring{$K3\times E$}{K3xE}]
\label{prop:k3e-separated-invariants}
The compact total space and the adjacent chiral, Borcherds, and fibre
constructions carry the following distinct invariants:
\[
\begin{array}{c|c|c}
\text{construction} & \text{invariant} & \text{value}\\
\hline
\text{compact total space} &
\kappa_{\mathrm{cat}}(K3\times E)=\chi(\cO_{K3\times E}) & 0\\
\text{compact Hodge supertrace} &
\kappa_{\mathrm{ch}}(K3\times E) & 0\\
\text{Heisenberg-Mukai chiral specialization} &
\kappa_{\mathrm{ch}}^{\mathrm{Heis}}(K3\times E) & 3\\
\text{Gritsenko--Borcherds denominator} &
\kappa_{\mathrm{BKM}}(\mathfrak{g}_{\Delta_5})=c_1(0)/2 & 5\\
\text{K3 fibre Mukai lattice} &
\kappa_{\mathrm{fiber}} & 24
\end{array}
\]
The value $2=\chi(\cO_{K3})$ belongs to the K3 fibre Euler witness; the
compact total-space invariant of $K3\times E$ is zero.
\end{proposition}

\begin{proof}
The total-space value is the K\"unneth computation in
Proposition~\ref{prop:k3e-current-ce}. For the compact Hodge supertrace,
Serre duality in odd complex dimension pairs $h^{0,q}$ with
$h^{0,3-q}$ and opposite signs, giving zero. The Heisenberg-Mukai
specialization is the framed $d=3$ chiral output associated to the K3
Mukai sector after elliptic specialization; its generator count gives
$3$. The Borcherds value is the weight of Gritsenko's
$\Delta_5$, equivalently $c_1(0)/2=10/2=5$ in the standard
normalization. The fibre value is the rank of the K3 Mukai lattice,
$24$.
\end{proof}

\section{Dimension and categorical scope}

\begin{proposition}[$d=3$ output scope]\label{prop:d3-output-scope}
For a CY$_3$ input equipped with the required orientation, framing,
completion, equivariance, and comparison data, the specialized
CY-to-chiral output is an $E_1$-chiral algebra. The $E_2$ structures
arise on Drinfeld centres, doubles, module categories, or bulk layers
under their own hypotheses.
\end{proposition}

\begin{proof}
The specialization from a holomorphic factorization algebra on the
CY$_3$ space to a chiral algebra on a curve uses a framed
two-dimensional transverse datum. The resulting algebra has ordered
one-dimensional chiral multiplication. Braiding requires the additional
centre or double construction, which supplies the half-braiding on the
representation category or the completed Hopf pairing on the double.
\end{proof}

\begin{remark}[Two proof lanes]\label{rem:two-proof-lanes}
The chain-level complex of Definition~\ref{def:chiral-ce-complex} and
the $(\infty,1)$-categorical Ran-space factorization statement of
Proposition~\ref{prop:ran-locality} are parallel load-bearing structures.
The first gives explicit differentials, filtrations, and supertraces.
The second gives descent, functoriality, and comparison with chiral
Koszul duality. A theorem using both lanes states the hypotheses in
each lane.
\end{remark}

\begin{thebibliography}{9}
\bibitem{BD04}
A.~Beilinson and V.~Drinfeld, \emph{Chiral Algebras},
AMS Colloquium Publications 51, 2004.

\bibitem{FG12}
J.~Francis and D.~Gaitsgory, \emph{Chiral Koszul duality},
Selecta Math.\ 18 (2012), 27--87.

\bibitem{CG17}
K.~Costello and O.~Gwilliam, \emph{Factorization Algebras in Quantum
Field Theory}, Vol.~1, Cambridge University Press, 2017.

\bibitem{Borcherds95}
R.~Borcherds, \emph{Automorphic forms on $O_{s+2,2}(\mathbb R)$ and
infinite products}, Invent.\ Math.\ 120 (1995), 161--213.

\bibitem{Gritsenko99}
V.~Gritsenko, \emph{Reflective modular forms and their applications},
Russian Math.\ Surveys 54 (1999), 639--640.
\end{thebibliography}

\end{document}
